\begin{center}
\LARGE
Regular Expressions Simplification

\end{center}

\begin{center}
\Large
Alejandro A.R. Trejo, Guillermo Fernandez Anaya

\end{center}

\noindent
Date:  July 18th (Thursday)\\
Time:  16:25-16:50\\

\begin{center}
\large
Abstract
\end{center}

Although Kleene's regular expressions are known since middle 50's and that play a central role in
many areas of computer science, such as in the design of sequential circuits and, very specially, in
the compiler theory, at the present time an authomatic method for their simplification is not
available. Such a method becomes desirable whenever you face involved language descriptions
which require long and tedius sequences of manipulations. 

This talk is devoted to present a technique to simplify regular expressions. Several ideas of
Computer Algebra and symbolic computing are applied to the regular expression symbolic
domain. An implementation of the described technique has been developed. 

The talk is divided in the following sections:
In section 1, the problem is stated. We present a wide view of the question proposed for solution. 
In section 2, the technique for simplifying is outlined. In section 3, the method of representation
and manipulation of expressions is explained. Finally, in section 4, the implemented regular
expression simplifier system is presented, together with examples of its usage. 

